\section*{Response to the September 26, 2026 R58 Report}
We thank the referee for the detailed mathematical, computational and repository audit. We agree that correctness and reproducibility are necessary but do not by themselves establish an Operations Research contribution. This revision addresses that distinction through a new original-model parameterized lower bound, a broader recognizable path-allocation class, optimized prototype selection, a direct convex comparator and a reorganized article. It does not withdraw the existing representation, compression, approximation, lattice or price-path results. Complete earlier readers and source files are preserved.

\paragraph{Governing version.}
The report is \path{operation_research_referee_report_r58_independent_harsh_2026-09-26.md} on the R58 independent-review branch, commit \nolinkurl{fe068d2a978a44548688b6052d4f210473394417}. The new author branch starts from the reviewed scientific tip \nolinkurl{1edd6d929e20f50a14ec6ef9beab4a37fa44fe39}. It does not merge reviewer-authored output into the author delta. The content map records where every current and retained scientific component is read.

\section{Scientific importance and the theory route: Sections 4 and 8.1}
We choose the theory route explicitly. Section~\ref{sec:class59} defines resource allocation on an arbitrary finite acyclic graph with endogenous source and path. The graph need not be a totally ordered command catalog, and the abstract problem contains no customer histories or prefix-eligibility assumption. Theorem~\ref{thm:recognition59} characterizes path-capacity conservation, supplies a linear-time recognition procedure and produces two different-capacity paths as a failure witness. It also identifies the additional source-balance condition needed for a common terminal deficit. Theorem~\ref{thm:abstractapprox59} gives an additive approximation and original-equality repair on this class, including a slower nonconcave Lipschitz version.

The renewal-specific theorem remains essential: it proves that the individual expected caps, realized ceilings and pre-draw service problem really has those edge payoffs and capacities. A general path identity alone would not provide that implementation. The article now states which assumptions are needed for the abstract conservation result and which are needed to derive the renewal kernels. We retain the service-credit architecture as a motivation, not as invented field calibration. The broader theorem and the original-model complexity result, rather than a claim of adoption, support disciplinary relevance.

\section{Limited-menu complexity: Sections 3.1, 3.5, 3.6 and 8.3}
Theorem~\ref{thm:wone59} is a new answer to the limited-menu question. It proves W[1]-hardness jointly in the command budget and the number of distinct expected caps, with those two parameters equal. The reduction remains within the original model and retains a common linear reward, zero service costs, uniform positive probabilities, ceilings equal to caps and nonnegative rational charges.

The key construction differs from the earlier nonbinding-budget reduction. A dedicated positive-probability history at each baseline forces that baseline into every threshold-achieving actual book. A zero-cap history forces the zero anchor. Equality in the value bound also allows at most one paid option per group. A no-carry multicolored-clique encoding then requires exactly one option from every group, so baseline commands cannot be added for free outside the cardinality accounting. Both the group encoding and the renewal mapping are proved in the main article.

Consequently the existing XP upper bound is not being offered as an unqualified tractability claim: an exponent independent of both parameters would imply FPT equals W[1]. We do not claim W[1] membership, a sharp exponent lower bound, strong NP-hardness, or a uniform-charge lower bound. Those distinctions are stated next to the theorem. The earlier weak SUBSET SUM construction, exact common-cap theorem and tight twice-cap-count compression remain in the companion. The new reduction is tested on 4,240 books and 48 graph instances; the general proof, not those finite checks, establishes the result.

\section{Prototype choice and approximation: Sections 3.2 and 9.4}
Section~\ref{sec:grouping59} makes prototype choice an explicit optimization problem. Catalog oscillation defines a pseudometric modulo additive constants, and cap-preserving representative selection has a standard medoid integer formulation. This makes the computational burden of general grouping visible instead of assuming free prototypes.

For quadratic rewards with a common curvature within each cap class, the oscillation cost reduces to a weighted absolute slope distance multiplied by catalog span. Weighted-median interval costs and a two-level dynamic program compute the exact distortion frontier over all group budgets in polynomial arithmetic time. The proof states the declared reward family and does not claim to solve arbitrary concave-function clustering. Combining its distortion with the existing approximate-type theorem gives an original-feasible menu-loss certificate and a bound on the required command union. All 160 tested frontiers agree with exhaustive medoid selection.

\section{Literature and contribution hierarchy: Sections 5, 7, 8.2 and 8.6}
The introduction now compares constrained shortest paths, separable resource allocation, moment-support theorems, auction menu complexity, choice-based assortment optimization, distributionally robust optimization, clustering and parameterized complexity. Table~\ref{tab:antecedents59} distinguishes established machinery from the conclusions proved here. New references are used to compare assumptions and theorem objects, not merely to enlarge the bibliography.

The introduction explicitly acknowledges that resource-state dynamic programming, Lagrangian branching and metric medoid formulations are antecedent tools. Paid unions of history-dependent terminal supports are distinguished from the support of one moment-constrained distribution and from incentive-compatible outcome--payment menus. Controlled command probabilities are distinguished from choice probabilities in assortment models. The payoff-oscillation guarantee does not purport to cover uncertainty in caps, probabilities or eligibility. Removing all artifact engineering still leaves the constructive reduction, the new parameterized theorem and the abstract recognition and approximation results.

The title, abstract and narrative now follow that spine. Essential new proofs remain in the main article. Secondary exact regimes, type refinements and complete price-search mathematics move to the companion; historical experiments and operational interpretation remain in the complete retained readers. No historical scientific file is silently deleted or overwritten. The final build records the actual page counts rather than assigning a submission category from an anticipated length.

\section{Computational redesign: Sections 6.1--6.3, 6.5--6.6 and 8.4}
The new panel uses 72 instances and 288 method requests. Eight seeds per principal cell replace the former three-seed design. Nonuniform rational probabilities, catalogs up to 65, histories up to 96, budgets up to eight, three positive accuracy levels and an exact-target group-sum family vary several relevant dimensions. Every input includes its distinct cap and joint-type counts, probability encoding and grid-size prediction. Inputs and executable sources are frozen before the new timed runs; the old study is not relabeled.

All competing methods receive the same three-second end-to-end allowance, including process startup, construction, serialization and checking. A 1.8-second internal optimization cutoff reserves time for evidence production; exceeding the total allowance is a failure to deliver a completed decision regardless of the last optimizer status. A direct mixed-integer convex model replaces the envelope-only comparator: Proposition~\ref{prop:micp59} proves exact equivalence and the code uses the quadratic epigraph directly. Its floating-point upper bounds remain explicitly numerical. The selected book is reevaluated with exact arithmetic and its original-space lower policy is independently checked. Thus formulation error is zero without claiming that a numerical solver status is an exact rational proof.

The primary table does not mix these evidentiary levels in one success column. The earlier equal-book-count lexicographic enumeration comparison is absent from the new main table; it remains in the unaltered historical reader. An exact structured oracle is available for the new group-sum challenge, but is not provided to timed methods. Memory, proof size, failed checks, state-limit rejections and deadlines are retained at row level. Checker RSS is clearly labeled as same-process high-water memory including retained optimizer allocations, not isolated incremental memory.

The new panel contains 288 method requests. Price search branches in 7 of 72 cases, with maximum recorded depth 16. Table~\ref{tab:study59} reports completed outputs and target attainment rather than optimizer status alone. Table~\ref{tab:checking59} retains end-to-end deadlines, grid rejections, and checking costs. The heterogeneous grid and binary-encoded challenge cases intentionally expose numerical-state and proof-cost limits. A common time allowance does not turn a numerical global bound into a rational certificate. These results identify method-dependent tradeoffs on this generator; they do not establish universal solver superiority or field performance.


This panel is a stronger controlled comparison, not a claim that three seconds is sufficient for every application or that synthetic instances are representative field data. The source freeze prevents timing relabeling, not all development adaptation. We do not call the study third-party or blind validation. The lower-bound theorem supplies the general complexity conclusion; observed branching depth is reported as evidence about these instances only.

\section{Screening: Sections 6.4, 8.5 and 9.9}
We retain the result that screening improved end-to-end time in none of the 18 earlier matched pairs and had a median ratio of approximately 1.604. The new main text calls this a negative preprocessing-economics result. Screening is not enabled in the new primary panel and is not presented as an empirically validated speedup. The safe fixing theorem and its original evidence remain available because correctness of a reduction and economy of using it are different scientific questions. We do not manufacture a favorable regime by selecting new outcomes.

\section{Specific technical and presentation comments}
\paragraph{Infeasible instances (comment 1).} Section~\ref{sec:infeasible59} states the necessary and sufficient anchor test before incumbent initialization. The new preflight implementation returns a digest-bound failed inequality, with no invented lower policy; malformed inputs are input errors. A separate checker recomputes the inequality. The existing solver exceptions remain documented for backward compatibility.
\paragraph{Complexity terminology (comments 2, 3 and 5--6).} The article separates polynomial per-price work, XP enumeration, the new parameterized lower bound and numerical-grid pseudo-polynomial complexity. Cap and joint-type counts appear in every input. The saturated regime states that positive probabilities and a saturated promise force every individual cap; lattice inputs expose both the numerical denominator and its bit length.
\paragraph{Comparison and checking (comments 7--10).} Equal-book-count enumeration is no longer a central comparator. The new direct quadratic model has zero envelope error and its solver gap is retained separately. Completed rational checking is required for rational target attainment. Conditional checking time and high-water memory are reported without hiding timeouts.
\paragraph{Related work, abstract and title (comments 11--13).} The related-work discussion and theorem-level comparison are in the article. The text-only abstract has 165 words and concentrates on one problem, one representation/complexity frontier and its algorithmic consequence. The title no longer suggests that small menus imply a generally efficient exact algorithm.
\paragraph{Units (comment 14).} The model uses a normalized obligation unit and a common net-value unit for reward, service cost and opening charges. An additive error is measured in that latter unit; multiplying all payoffs and the tolerance by the same positive scale leaves the decision guarantee invariant. No currency calibration is asserted without data.

\section{Preservation and exact-head qualification: Sections 1.3, 8.7 and comment 15}
The release retains exact predecessor copies of every changed root reader and keeps all prior scientific paths unchanged. The code-and-data archive opens directly onto the current root readers and their dependency paths; a referee need not discover a nested archive. Build and source manifests bind the compiled reader graph, mathematical labels and measured diagnostics.

The production step publishes a candidate commit. A separate read-only verification step checks that exact commit's frozen sources, complete records, input hashes, proof bytes, reader hashes and preservation manifest, and attaches a commit status to that SHA. It does not describe a build of an earlier head as qualification of the final head. Making that context administratively required by branch protection is a repository-owner setting; the publication procedure does not claim to have changed protection policy. Neither a passing check nor this response claims editorial acceptance.
